% =====================================================================
% Plano de desenvolvimento --- Vitória (Frente 2)
% Só a Vitória edita este arquivo. A tabela de marcos foi copiada da divisão
% de tarefas; a Vitória revisa e completa com o plano dela.
% =====================================================================
\section{Plano de desenvolvimento}\label{sec:plano}

O trabalho é dividido em seis frentes de código, uma por integrante, que avançam em paralelo
a partir da classe \codigo{Grafo}. Cada função entra na \codigo{main} por \emph{pull request},
revisada por outro integrante e com os testes passando. O andamento é acompanhado num
cronograma compartilhado. A Tabela~\ref{tab:marcos} mostra os marcos.

\begin{table}[htbp]
\centering
\caption{Marcos do trabalho. A entrega final no Canvas é até 22/11/2026.}
\label{tab:marcos}
\small
\begin{tabularx}{\textwidth}{@{}l l Y@{}}
\toprule
Marco & Data & O que precisa estar pronto \\
\midrule
Base pronta (Etapa 1) & 10/11 & Classe \codigo{Grafo} e leitores na \codigo{main}; este documento \\
Algoritmos (Etapa 2) & 17/11 & Funções de cada frente na \codigo{main}, com os testes batendo com os valores de referência \\
Integração (Etapa 2) & 19/11 & Interface com as seis telas; SIR e ranking na Twitch PT; medições de desempenho \\
Entrega (Etapa 3) & 21/11 & Relatório SBC de 12 a 15 páginas, slides, vídeo e o \codigo{.zip} no Canvas \\
\bottomrule
\end{tabularx}
\end{table}

\paragraph{Fases e dependências.} Na primeira fase, a Frente 1 entrega a classe \codigo{Grafo}, os
leitores e os geradores, de que todas as outras dependem; enquanto isso, as demais frentes estudam
os algoritmos e preparam os testes com os valores de referência. Na segunda, as Frentes 2 a 4 avançam
em paralelo, porque só dependem do \codigo{Grafo} e de um contrato comum: cada função recebe um
\codigo{Grafo}, não o altera e devolve um valor por nó (\codigo{dict[nó, valor]}). A Frente 5 depende
das seis centralidades das Frentes 2 e 3, por isso começa pelo SIR e pelo $\tau_b$, e a Frente 6 monta
as telas com dados de exemplo e passa a usar as funções reais à medida que elas entram na
\codigo{main}. Na terceira fase, a interface integra todas as funções, a Frente 3 mede o desempenho
contra o NetworkX e a Frente 5 roda o experimento na Twitch PT. Na última, cada integrante escreve as
suas seções do relatório final e grava a sua parte do vídeo.

\paragraph{Riscos.} (1) O Python puro pode ser lento demais na intermediação exata, de custo $O(nm)$:
a função aceita sortear só $k$ origens. (2) Um erro numa frente pode contaminar as outras: cada
\emph{pull request} precisa de testes que comparem o resultado com valores conhecidos ou com o
NetworkX. (3) O atraso de uma frente pode travar a interface: o contrato comum permite usar dados de
exemplo até a função ficar pronta. (4) O prazo é curto: o cronograma compartilhado mostra o que está
atrasado e o que bloqueia os outros.
